\chapter{Variational Autoencoders}

\section{Introduction}

Variational Autoencoders (VAEs) combine latent-variable generative modeling with neural networks and variational inference. They provide a principled framework for learning probabilistic generative models while remaining computationally tractable.

This chapter develops VAEs from first principles: we derive the Evidence Lower Bound (ELBO), introduce the encoder--decoder architecture, and explain the reparameterization trick. We also interpret the tradeoff between reconstruction and regularization and discuss common failure modes.

\section{Autoencoding and Latent-Variable Modeling}

A latent-variable generative model assumes:
\[
p_\theta(x) = \int p_\theta(x|z)p(z)\,dz.
\]

In VAEs:
\begin{itemize}
    \item $p(z)$ is a simple prior (e.g., $\mathcal{N}(0,I)$),
    \item $p_\theta(x|z)$ is a neural decoder.
\end{itemize}

\subsection{Autoencoding View}

We introduce an encoder $q_\phi(z|x)$ that maps data to latent variables. Thus:
\begin{itemize}
    \item encoder: $x \to z$,
    \item decoder: $z \to x$.
\end{itemize}

Unlike classical autoencoders, this mapping is probabilistic.

\section{Variational Inference and the ELBO}

The log-likelihood is:
\[
\log p_\theta(x) = \log \int p_\theta(x,z)\,dz.
\]

Introduce an approximate posterior $q_\phi(z|x)$:
\[
\log p_\theta(x)
=
\log \int q_\phi(z|x)\frac{p_\theta(x,z)}{q_\phi(z|x)}\,dz.
\]

Applying Jensen's inequality:
\[
\log p_\theta(x)
\ge
\mathbb{E}_{q_\phi(z|x)}\left[\log \frac{p_\theta(x,z)}{q_\phi(z|x)}\right].
\]

\begin{definition}[ELBO]
\[
\mathcal{L}(x;\theta,\phi)
=
\mathbb{E}_{q_\phi(z|x)}[\log p_\theta(x|z)]
-
\mathrm{KL}(q_\phi(z|x)\|p(z)).
\]
\end{definition}

\subsection{KL Decomposition}

We have:
\[
\log p_\theta(x)
=
\mathcal{L}(x;\theta,\phi)
+
\mathrm{KL}(q_\phi(z|x)\|p_\theta(z|x)).
\]

Thus maximizing the ELBO both:
\begin{itemize}
    \item increases likelihood,
    \item improves approximation of the true posterior.
\end{itemize}

\section{Encoder--Decoder Architecture}

\begin{itemize}
    \item Encoder: $q_\phi(z|x)$ outputs parameters of a distribution over $z$,
    \item Decoder: $p_\theta(x|z)$ generates data from latent variables.
\end{itemize}

\subsection{Typical Gaussian Encoder}

\[
q_\phi(z|x) = \mathcal{N}(z \mid \mu_\phi(x), \Sigma_\phi(x)).
\]

\subsection{Decoder}

\[
p_\theta(x|z) = \text{neural network output distribution}.
\]

Thus learning involves jointly optimizing $\theta$ and $\phi$.

\section{Reparameterization Trick}

Direct gradient estimation of the ELBO is difficult due to sampling from $q_\phi(z|x)$.

\subsection{Reparameterization}

For Gaussian latent variables:
\[
z \sim \mathcal{N}(\mu, \sigma^2)
\quad \Longleftrightarrow \quad
z = \mu + \sigma \odot \epsilon,\quad \epsilon \sim \mathcal{N}(0,I).
\]

\subsection{Reparameterization Estimator}

The expectation becomes:
\[
\mathbb{E}_{q_\phi(z|x)}[f(z)]
=
\mathbb{E}_{\epsilon \sim \mathcal{N}(0,I)}[f(\mu_\phi(x) + \sigma_\phi(x)\epsilon)].
\]

This allows gradients to flow through $\mu_\phi$ and $\sigma_\phi$.

\section{Full VAE Objective}

The ELBO becomes:
\[
\mathcal{L}(x)
=
\mathbb{E}_{\epsilon}\left[\log p_\theta(x|z)\right]
-
\mathrm{KL}(q_\phi(z|x)\|p(z)).
\]

\subsection{Interpretation}

\begin{itemize}
    \item first term: reconstruction accuracy,
    \item second term: regularization toward prior.
\end{itemize}

Thus VAE balances:
\[
\text{reconstruction} \quad \text{vs} \quad \text{latent regularization}.
\]

\section{Gaussian Latent VAE}

For:
\[
q_\phi(z|x)=\mathcal{N}(\mu,\sigma^2),
\quad
p(z)=\mathcal{N}(0,I),
\]
the KL term has closed form:
\[
\mathrm{KL}
=
\frac{1}{2}\sum_j \left(
\mu_j^2 + \sigma_j^2 - \log \sigma_j^2 - 1
\right).
\]

This makes training efficient and stable.

\section{Interpretation of the KL Term}

The KL term:
\begin{itemize}
    \item encourages $q_\phi(z|x)$ to stay close to $p(z)$,
    \item prevents overfitting,
    \item enforces a structured latent space.
\end{itemize}

\subsection{Tradeoff}

\begin{itemize}
    \item strong KL: better regularization, worse reconstruction,
    \item weak KL: better reconstruction, risk of overfitting.
\end{itemize}

\section{Failure Modes and Variants}

\subsection{Posterior Collapse}

The encoder ignores $x$ and outputs:
\[
q_\phi(z|x) \approx p(z).
\]

Then:
\begin{itemize}
    \item latent variables carry little information,
    \item model behaves like an unconditional generator.
\end{itemize}

\subsection{Causes}

\begin{itemize}
    \item overly powerful decoder,
    \item strong KL regularization.
\end{itemize}

\subsection{Variants}

\begin{itemize}
    \item $\beta$-VAE: scale KL term,
    \item KL annealing,
    \item richer posterior families.
\end{itemize}

\section{A Unifying Perspective}

VAEs combine:
\begin{itemize}
    \item latent-variable generative modeling,
    \item variational inference,
    \item neural parameterization.
\end{itemize}

The ELBO provides a tractable objective balancing data fit and latent structure. The reparameterization trick enables efficient gradient-based learning.

Thus VAEs bridge probabilistic modeling and deep learning.

\section{Summary}

In this chapter we developed variational autoencoders. We derived the ELBO, introduced encoder--decoder architectures, and explained the reparameterization trick. We interpreted the KL term and the reconstruction tradeoff, and discussed failure modes such as posterior collapse.

VAEs provide a principled and flexible framework for generative modeling with latent variables.